\section{From Survey Measures to an Inflation Sentinel: Capacity Utilization during the Fordist Era}

Making idle productive capacity visible preceded its stabilization as a macroeconomic variable. Early efforts, such as the \citet{Nourse1934} Brookings Institution studies, immediately revealed a tension between engineering definitions of maximum throughput and economic notions of sustainable operation \citep{Burns1954; KleinDavid1958}. As postwar stabilization policy required a metric to translate idle capacity into intervention thresholds, capacity utilization filled that role. However, the measure that became institutionalized was not anchored in physical reality.

\citet{Butler1958} and \citet{Phillips1963} note that McGraw--Hill's plant-and-equipment surveys regularized managerial reports of ``preferred'' operating rates, translating heterogeneous capacity definitions into a standardized reporting stream. Crucially, this institutionalization was not a neutral technical upgrade. McGraw--Hill operated as a central node in the informational infrastructure of monopoly capital, leveraging its expanding business intelligence networks to anchor the measure’s legitimacy in corporate forecasting and policy debate \citep{BaranSweezy1988; Munroe2007}. By sedimenting these administrative routines, the survey marginalized deeper methodological debates and physically anchored alternatives, effectively naturalizing subjective managerial expectations as the de facto statistical fact of economic slack.

The Federal Reserve subsequently embedded these survey-based benchmarks into the Industrial Production program \citep{Board2017}, a choice reflecting administrative path dependence and institutional reach rather than empirical superiority. Policymakers privileged capacity utilization over the unemployment rate because it provided a real-time diagnostic of capital bottlenecks and production constraints, whereas labor market indicators were viewed as lagging social outcomes. Physically anchored alternatives---such as \citet{Foss1963}'s Electric Motors Utilization or \citet{TaubmanGottschalk1971}'s Average Workweek of Capital---tracked operational realities more directly but were systematically marginalized. 

The fiscal overheating of the late 1960s and the subsequent stagflationary crisis of the 1970s did not dismantle this survey routine; they recoded its policy function. Amid these compounding crises, the indicator's primary function was recast: it ceased to be a gauge of idle resources and became the central predictor of inflationary pressure \citep{KleinEtAl1973}. This durability did not resolve the underlying identification problem. Instead, it cleared institutional space for a new governance architecture---one that would displace survey-centered utilization with output-gap benchmarks portable across central banks, a transition detailed in the following section.

\section{From Capacity Utilization to Output-Gap Governance}

The erosion of Bretton Woods capital controls dissolved the territorial anchor of Fordist measurement. As the institutional apparatus of discretionary demand management was politically dismantled, giving way to rules-based frameworks, macroeconomic statecraft required a new benchmark for economic slack---one capable of circulating across central banks and supranational institutions independently of national industrial surveys \citep{Jessop2002; Jessop2003; Konings2007}. This displacement was not a technical upgrade, but a reorganization of governance around a \textit{portable architecture of slack}. 

The epistemic shift was crystallized in the rational expectations revolution, which redefined the output gap not as a measure of idle resources, but as a deviation from a latent, model-implied natural rate \citep{Johnson2024}. This was not a clean theoretical rupture, but an instance of \textit{institutional conversion} \citep{MahoneyThelen2009}: the same optimization-stability toolkit that \citet{Samuelson1947} had championed to legitimize Keynesian demand management was strategically redeployed to enforce credibility-based discipline. The rhetorical reorientation was captured perfectly by Alan Greenspan at a 1973 Brookings symposium, who argued that if a capacity measure did not signal inflationary pressure, ``then I suggest that the numbers are wrong'' \citep[p.~758]{KleinEtAl1973}. Measurement was no longer tasked with describing reality; it was tasked with enforcing a market-clearing benchmark.

The output gap consolidated not because it solved the measurement problem, but because it could be produced through multiple, mutually reinforcing closure strategies. Actual output ($Y$) is recorded; potential output ($Y^*$) must be inferred. The output gap did not merely track $Y$; it operationalized the \textit{distance} between $Y$ and an unobservable $Y^*$, creating a portable architecture of slack. Whether through statistical filtering \citep{HodrickPrescott1997}, neoclassical production functions \citep{CBO2004PotentialGDP}, or structural VARs \citep{BlanchardQuah1989}, each method imposes closure assumptions that mechanically dictate the permissible distance between actual activity and capacity. Methodological divergence masks functional convergence: each route collapses the politics of capacity identification into a single, portable variable.

The Global Financial Crisis exposed the empirical fragility of this settlement. Post-crisis revisions routinely flipped the sign of the output gap, inverting the policy signals the benchmark was designed to deliver \citep{CerraFatasSaxena2023; BarkemaGudmundssonMrkaic2020}. Patching the apparatus required discretionary ``nowcasting'' interventions \citep{AastveitTrovik2014; Berger2023} that steadily eroded the claim to rule-bound expertise on which its depoliticized authority rested \citep{Foster2024; ArgyroulisVagdoutis2025}. The durability of these practices signals institutional entrenchment, not empirical adequacy. It is within this closure-dependent, revision-prone setting that political economy's search for a structurally identified, accounting-based measure of capacity becomes historically intelligible.

\section{Capacity-Utilization Debates in Political Economy}

The utilization controversy divides along a structural fault line: whether the normal rate of capacity utilization is endogenous to demand or anchored in cost structures. The Neo-Kaleckian tradition treats the normal rate as demand-determined, where accumulation trajectories lock into path-dependent outcomes \citep{Lavoie1995; Nikiforos2013; Setterfield2020}. Conversely, the Sraffian-Keynesian tradition anchors the normal rate in production costs, positing that actual utilization gravitates toward this center over the long run, though not without facing its own closure problems regarding autonomous demand \citep{Nikiforos2018}. 

Both traditions face a shared empirical impasse regarding the long-run business cycle: resolving this theoretical split requires a reliable measure of long-run utilization, but available mainstream proxies are fundamentally ill-equipped for the task. The Federal Reserve series cannot supply this instrument because its statistical behavior reflects administrative design, not underlying economic dynamics. As \citet{Shapiro1989} demonstrate, the FRB construction process blends survey data on current operating rates with regression-based capacity estimates smoothed against capital stock and time trends, explicitly calibrated to approximate sustainable output under normal conditions. 

This methodological design hardwires mean-reversion into the series, rendering it incapable of registering chronic departures from normal utilization. Consequently, testing the stationarity of the FRB series does not verify economic convergence; it merely verifies the Federal Reserve’s own smoothing protocol \citep{Nikiforos2016; Nikiforos2021}. The controversy requires evidence across capital-depreciation horizons, yet the series operates as a cyclical slack indicator optimized for short-run policy surveillance. 

While physical indicators like the average workweek of capital track machinery use more directly and exhibit trending behavior where the FRB series remains stationary by construction, single-series strategies still fail to recover a structural relation jointly shaped by utilization, demand, and accumulation. If the long-run relationship is a system property, it cannot be identified through isolated proxies. This proxy impasse necessitates a structural measurement route that bypasses institutional surveys entirely, shifting the burden of capacity estimation from administrative routines to profitability accounting.

\section{Shaikh’s Structural Identification of Capacity Utilization}

To resolve this proxy impasse, \citet{Shaikh2016} shifts capacity estimation into profitability accounting. Rather than relying on institutional surveys or unobservable potential output, he reconstructs the capital stock through a generalized perpetual inventory method (GPIM) and derives a theoretical capacity path from the accounting link between actual and maximum profit rates at the normal cost of capital. Utilization is then constructed as the residual deviation from this implied output–capital ratio.

This accounting route gains immediate analytical traction through data portability. Because it relies on standard national accounts for output and capital stock, it bypasses institutional survey constraints and applies consistently across jurisdictions. This portability enables robust series construction where physical proxies are unavailable, while embedding the theoretical assumption that normal utilization adjusts to demand conditions rather than remaining fixed by technology.

However, the residual construction establishes an accounting identity but leaves identification open. Defining utilization as a deterministic residual imposes a capacity path without verifying whether the underlying variables actually share a stochastic equilibrium. When identification remains weak, the single-equation estimation becomes unstable, risking the registration of constructed residuals as false cointegration signals. The long-run output–capital coefficient is forced to absorb shifting distributive regimes, generating a highly fragile empirical benchmark.

This article is motivated to critically engage with \citet{Shaikh2016}'s cointegration approach to stress-test this fragility. By reinterpreting the long-run output–capital coefficient as a transformation elasticity ($\theta$) linking capital accumulation to productive-capacity formation, this paper evaluates whether a fixed-parameter closure can sustain structural meaning across postwar distributive regimes. The multi-stage replication developed in the following sections demonstrates that the standalone output–capital system fractures due to severe omitted variable bias, proving that valid capacity identification must place distributional variables—specifically, the logged rate of exploitation—directly inside the structural estimation architecture.